\section{Introduction}

Vision--language--action (VLA) policies map visual observations and language
instructions to action chunks, inheriting semantic and control priors from large
pretrained models.  World models address a complementary problem: they predict
how actions change future state and enable policy optimization through imagined
trajectories.  Combining them requires a precise boundary between the policy's
internal decoder state and the observation state learned by the world model.

That boundary is often obscured by checkpoint-specific hidden tensors, flattened
features, or multiple preprocessing routes.  Such ambiguity is especially risky
in online cotraining, where collectors, replay buffers, world-model learners,
success classifiers, actors, and evaluation workers must agree on shape and
semantics.  A tensor that is convenient for decoding actions is not necessarily
a valid observation of the environment.

DreamerVLA uses the projected visual input tokens of a one-trajectory
OpenVLA-OFT checkpoint as its sole external observation representation.  The
policy consumes one current agent-view image and produces a $256\times4096$
token grid before its internal action positions.  Reward and token sidecars are
paired exactly by shard, demonstration, and frame.  A chunk-aware world model
predicts future grids, a spatial classifier estimates success on windows of the
same grids, and an actor bridge decodes predicted grids through OpenVLA's
discrete action vocabulary.

The complete training path is
\emph{collect rollouts $\rightarrow$ seed replay $\rightarrow$ warm up the world
model and classifier $\rightarrow$ online cotrain}.  Synchronous and asynchronous
executions share the same update functions and checkpoint contracts.  The
asynchronous route assigns policy training, no-gradient inference, environment
interaction, and world-model/classifier learning to explicit worker groups.
